\section{Exploratory Data Analysis}

The EDA was structured into four passes, mirroring the natural questions a
data miner asks when meeting a dataset for the first time: structural,
distributional, content-metadata, and temporal. Full reports and 26 plots
live in \texttt{output/eda/}; this section consolidates the most
load-bearing findings.

\subsection{Structural overview}

Of the 8 columns, three carry the bulk of the missing-value mass —
\texttt{cursor\_time} (71.6\% missing), \texttt{user\_answer} (78.5\%), and
\texttt{source}/\texttt{platform} (0.02\% each). Crucially, all of these are
\textbf{structurally missing}: \texttt{cursor\_time} is only emitted by media
actions, \texttt{user\_answer} only by respond/erase actions, and
\texttt{source}/\texttt{platform} are null only on payment-related actions
(no learning context). Within their applicable subsets they are 100\% present.
This means \textbf{no imputation is needed} — filling these would inject
misleading information rather than recover lost signal.

\subsection{Distributions and the long-tail of users}

User activity is extremely right-skewed: the median user has only \textbf{31
interactions}, while the mean is \textbf{441}. The top 1\% of users contribute
8{,}300+ interactions each, dragging the mean up by an order of magnitude.
This has direct consequences for downstream modeling: cold-start users
(P25 = 22 interactions) carry very little behavioral history, while a
small power-user cohort dominates any naive aggregate statistic.

\begin{figure}[H]
\centering
\includegraphics[width=0.62\linewidth]{user_activity.png}
\caption{User-level interaction count distribution (log-y). The long tail makes cold-start vs.\ power-user a first-class modeling concern.}
\end{figure}

Other distributional findings worth flagging:

\begin{itemize}
\setlength\itemsep{1pt}
\item \textbf{Sprint mode dominates} (71.3\% of all actions). Students
  prefer self-directed part-specific practice; the AI's
  \texttt{adaptive\_offer} accounts for only 6.6\%, and \texttt{tutor} just
  0.9\%. Most learning is student-driven, not platform-driven.
\item \textbf{Mobile/web split: 71/29}, stable across nearly every
  action type. The single exception is \texttt{undo\_erase\_choice} —
  142{,}092 events on web, \textbf{exactly zero on mobile} — indicating
  the feature was never shipped on the mobile UI.
\item \textbf{Audio dominates over video} (~25\% of all actions vs.\ ~3\%),
  consistent with TOEIC's listening-heavy format.
\end{itemize}

\subsection{Content metadata: questions, parts, difficulty}

The 13{,}169 questions are heavily concentrated in \textbf{Part 5
(Incomplete Sentences)}, which holds 43.3\% of the bank — nearly half of
all questions are grammar/vocabulary fill-in-the-blank. Listening parts
1--4 are smaller individually but use bundled passages (multiple questions
per audio clip).

Difficulty was estimated as $1 - \text{accuracy rate}$ across all student
responses. Across 23.3M matched responses, \textbf{overall accuracy is
56.87\%}; difficulty is approximately normal around 0.39, with very
few trivially easy or impossibly hard items.

\begin{table}[H]
\centering
\small
\begin{tabular}{llrr}
\toprule
Part & Section & Avg. difficulty & Accuracy \\
\midrule
5 & Incomplete Sentences   & \textbf{0.487} & 51.3\% \\
6 & Text Completion        & 0.404 & 59.6\% \\
7 & Reading Comprehension  & 0.393 & 60.7\% \\
4 & Short Talks            & 0.377 & 62.3\% \\
2 & Question-Response      & 0.369 & 63.1\% \\
3 & Short Conversations    & 0.364 & 63.7\% \\
1 & Photo Descriptions     & \textbf{0.326} & 67.4\% \\
\bottomrule
\end{tabular}
\caption{Difficulty by TOEIC part. Part 5 is both the largest and the hardest section.}
\end{table}

The combination is striking: Part 5 is simultaneously the
\emph{largest} share of the question bank (43\%) and the \emph{hardest}
to answer (51.3\% accuracy). Students spend the most time exactly where
they perform the worst.

\subsection{Temporal patterns}

Daily activity exhibits a clear bimodal pattern across the 461-day
collection window, with peaks in \textbf{January 2019} and
\textbf{July--August 2019}. These align with common TOEIC exam
preparation cycles in Korea — students ramp up usage in the months
preceding major test sittings. August 2019 is the busiest single month
at 14M interactions.

\begin{figure}[H]
\centering
\includegraphics[width=0.78\linewidth]{monthly_activity.png}
\caption{Monthly interaction volume. Bimodal peaks (Jan 2019, Jul--Aug 2019) align with TOEIC exam-prep seasons.}
\end{figure}

Day-of-week and hour-of-day analyses show \textbf{weekday domination}
(Tuesday is the busiest day, Sunday the quietest, ~23\% lower) and a
\textbf{single-peak daily cycle} — there is no secondary evening study
peak. Daily Active Users averages just 2{,}013 against the 297K total
user base (~0.7\% DAU/MAU), typical for a freemium edtech product with
high churn outside exam season.
